% !TeX root = ../../../geos_mpm_manual.tex
\section{Native XY plane strain}
\label{sec:native-true-2d}
\index{plane strain!native 2D}
\index{planeStrainThickness}
\index{CPTI!triangles}

\subsection{One timestep, dimension-specific kernels}
\texttt{formulation=Cartesian3D} remains the default. Selecting
\texttt{PlaneStrain} uses the same timestep sequence with dimension-specialized
mapping, gather/update, integration-domain and splitting paths. The background
mesh has native C2D4 cells and two-node faces. Partition ownership, periodic
neighbors and ghost regions operate in XY; no z cell or z communication slab
is constructed. Framework vectors retain three components with inactive z zero.
Constitutive tensors remain $3\times3$ so that plane-strain material response
retains $\sigma_{zz}$ while $F_{zz}=1$ and $L_{zz}=0$.

G and $G_0$ each store two full XY edge generators. Shape 3 reconstructs a
parallelogram at $c\pm g_1/2\pm g_2/2$. Shape 4 reconstructs a triangle at
$v_0=c-(g_1+g_2)/3$, $v_1=v_0+g_1$, $v_2=v_0+g_2$. No thickness generator,
vertex array or glyph array is stored. The two matrices cost 64 bytes per
particle. F-only and F-plus-geometry reset semantics are unchanged, applied
to the active geometry block. Native splitting retains shape and interpolation.

\begin{center}\small
\begin{tabular}{@{}lrrp{0.35\linewidth}@{}}
\toprule Mapping & 2D entries & 3D entries & Material geometry\\\midrule
SinglePoint & 4 & 8 & Any matching-dimension shape\\
SinglePointBSpline & 16 & 64 & Any matching-dimension shape\\
CPDI & 16 raw & 64 raw & Parallelogram / box\\
CPTI & 12 raw & 32 raw & Triangle / tetrahedron\\
\bottomrule
\end{tabular}\end{center}
Raw CPDI/CPTI entries include repeated grid nodes. Effective mapping coalesces
them. For triangle CPTI, $E=[g_1\ g_2]$, $b_a=E^{-T}e_a$ and
$b_0=-b_1-b_2$. With bilinear grid basis $N_I$, the mapping is
\[
 S_I=\frac{1}{3}\sum_{a=0}^{2}N_I(x_a),\qquad
 \nabla S_I=\sum_{a=0}^{2}N_I(x_a)b_a.
\]
Corners and gradient coefficients use the same temporary integration geometry.
CPDI uses four corners and the reciprocal parallelogram basis. Homel and
VP-Hencky retain their active-dimension box policies. CPTI uniformly clips
its diameter using active grid spacings or neighbor radius, without the box
volume-preserving feasibility cap. Scaling never changes stored material
geometry, physical volume, mass or F.

\subsection{Physical thickness and statistical volume}
The single finite positive \texttt{planeStrainThickness} is written on
\texttt{InternalMesh} and stored on the domain partition, including restart.
It supplies
\[
 V_p=A_p t,\qquad m_p=\rho_p A_p t,\qquad
 A_\Gamma=\ell_\Gamma t.
\]
In-plane geometry, CPDI scaling, diameter-based splitting and plot extents are
independent of $t$. The particle-sized fracture \texttt{lengthScale} uses
$\sqrt{V_p/t}$; a scalar constitutive length-scale parameter is not overwritten.
Nodal physical-volume normalization uses $h_xh_yt$. CZ line-interface weights
integrate the Q1 basis along the segment and multiply by $t$. Force and energy
retain physical units, while constitutive tensors include the out-of-plane
plane-strain stress.

Native dim=2 Voronoi/layered Weibull wrappers use $\max(A t,V_{\min})$ in the
strength law. They preserve random draws when the writer binds thickness;
rebinding is idempotent. An explicit old layered-wrapper thickness must agree
with the domain scalar. Arbitrary supplied \texttt{StrengthScale} fields are
not rescaled; custom fields can implement \texttt{setPlaneStrainThickness(t)}.
The Brazilian-disk and borehole examples now use native 2D random fields, so
their realization changes from the previous 3D construction even with the same
seed. General 3D stochastic/flaw-file models retain their previous meaning.

\subsection{Objects without artificial end caps}
New \texttt{geom.disk} represents a disk or annulus using an XY center and radii;
\texttt{geom.rectangle} constructs \texttt{box(dim=2)}. Existing
\texttt{polygon}, \texttt{foam(dim=2)}, and
\texttt{packedSphericalBed(dim=2)} already perform XY geometry queries.
Native flags, normals and surface offsets do not depend on query z.
\begin{lstlisting}[language=Python]
disk = geom.disk('disk', [0.0, 0.0], 1.0, particleType=2)
ring = geom.disk('ring', [0.0, 0.0], 1.0, ri=0.4,
                 flaggedSurfaces=[True, True]) # outer, inner
block = geom.rectangle('block', [-1.0, -0.5], [1.0, 0.5])
\end{lstlisting}
Before sampling PlaneStrain objects, the writer adapts simple legacy z-axis
cylinders to disk/annulus sections and axis-aligned 3D boxes to XY rectangles.
It removes end-cap influence from membership and surface methods and traverses
owned wrapper subobjects. This assumes these objects represent intended
extrusions. Material/group/velocity attributes remain on the original objects.
The disk frame matches the positive-z legacy cylinder frame.

Tilted cylinders, oblique caps, spheres, STL/general 3D objects, and transforms
mixing XY with z retain existing cross-section sampling. They require deliberate
section geometry; projecting an arbitrary 3D normal alone is not an exact
section algorithm. See \path{PFW_2D_OBJECTS.md} for the example audit and the
legacy adaptation contract. The copper-foam, packed-bed and polygon examples
retain their native XY objects; disk examples now use the new constructor.

\subsection{Input, output and restart contract}
Native inputs need no z bounds or nK. The example
\path{examples/true2D/pfw_input_true_2d.py} combines CPDI parallelogram bulk with
CPTI surface triangles using \texttt{simplex},
\mpmcode{particleGeometry=fromRefinement}, and \mpmcode{simplexParticleType=3}.
The named hexTet preset remains 3D only. All-CPTI objects tessellate each
Cartesian seed into two triangles in 2D or six tets in 3D.

Old Python \texttt{planeStrain=1} maps to the native formulation. If thickness
is unset, its old z extent supplies $t$; otherwise new inputs default to 1.
Four boundary types and time-plus-four boundary tables are supported; old
six-face values are accepted with z inactive. F tables retain four columns
(time and three diagonal entries), with Fzz=1. The optional transverse-velocity
matrix retains its existing 6-by-2 interface; only in-plane entries act.

Native particle rows carry version 2, shape 3/4, physical volume, and only four
geometry components XX/XY/YX/YY. PositionZ is zero or may be omitted by a native
reader. Both geometry matrices are restartable. Restart geometry version 3 is
required; old GEOS particle/XML/restart compatibility is not retained.
Regenerate inputs and start a new run after applying this patch.

Silo writes dimension-2 quad/triangle zones. VTK writes quad/triangle cells with
z=0 coordinates. Both reconstruct current unscaled G, and grid quads use
consistent perimeter order. The shared VisIt script selects a fitted
\texttt{View2D} from native-2D metadata; case renderers accept four-value XY
extents. Full material tensors, including $\sigma_{zz}$, remain available.

\subsection{Validation boundary}
Standalone CPU tests compile production geometry/mapping code and selected
mesh/partition/splitting bodies with small adapters. They check mapping
moments, compact storage, reset/scaling, native split conservation, line
integrals, quad topology, partition neighbors, object surface data and Weibull
volume scaling. Writer checks use explicitly emulated ranks, including empty
ranks. Camera tests mock the VisIt API.

These checks do not replace the target GEOS build, actual MPI migration and
restart, GPU execution, coupled contact/CZ mechanics, or
Silo/VTK readback in VisIt. The delivery validation report lists exact local
results. Axisymmetric single-field mechanics are described in
Section~\ref{sec:axisymmetric}; its remaining extensions, plane stress, CPDI2
and cohesive-particle subdivision remain separate work. See \path{TRUE_2D_PLAN.md} for target acceptance cases.
